\section{Tercera parte: construir un Hammer para Lean --- LeanHammer}

\subsection{Arquitectura del pipeline de Hammer}

\begin{importantbox}{Componentes de LeanHammer}
\begin{enumerate}
    \item \textbf{Selección de premisas} (Premise Selection): un modelo neuronal de recuperación aplica embedding al objetivo y a las premisas candidatas, y usa similitud coseno para seleccionar los teoremas, definiciones o lemas más relevantes
    \item \textbf{Capa de traducción} (lean-auto): traduce la teoría de tipos dependientes de Lean a la lógica del demostrador automático
    \item \textbf{Demostrador automático de teoremas} (Zipperposition, entre otros)
    \item \textbf{Reconstrucción de la demostración} (Duper): traduce la salida del demostrador automático de vuelta a una demostración en Lean
    \item \textbf{Búsqueda en árbol} (AESOP): organiza todo el proceso de búsqueda
\end{enumerate}
\end{importantbox}

El usuario solo necesita invocar el comando \texttt{hammer} dentro de la demostración, y el sistema realiza automáticamente la selección de premisas, la traducción, la demostración automática y la reconstrucción.

\subsection{Resultados experimentales}

\begin{itemize}
    \item El selector neuronal de premisas de LeanHammer (con solo 100 millones de parámetros) supera claramente a los trabajos anteriores
    \item Es complementario al selector simbólico de premisas (Moogle); al tomar la unión de ambos, el desempeño mejora aún más
    \item El desempeño se está acercando al límite teórico superior que se obtiene al usar las premisas reales presentes en demostraciones humanas
\end{itemize}

\subsection{Resumen de la sección}

DSP y LeanHammer representan dos enfoques complementarios: el primero guía la estructura de la demostración mediante bosquejos de alto nivel, mientras que el segundo ofrece una herramienta de automatización de bajo nivel, poderosa, para completar los vacíos.
